\documentclass[10pt]{article}  

%%%%%%%% PREÁMBULO %%%%%%%%%%%%
\title{Memoria para prácticas de DOMÓTICA}
\usepackage[spanish]{babel} %Indica que escribiermos en español
\usepackage[utf8]{inputenc} %Indica qué codificación se está usando ISO-8859-1(latin1)  o utf8  
\usepackage{amsmath} % Comandos extras para matemáticas (cajas para ecuaciones,
% etc)
\usepackage{amssymb} % Simbolos matematicos (por lo tanto)
\usepackage{graphicx} % Incluir imágenes en LaTeX
\usepackage{color} % Para colorear texto
\usepackage{subfigure} % subfiguras
\usepackage{float} %Podemos usar el especificador [H] en las figuras para que se
% queden donde queramos
\usepackage{capt-of} % Permite usar etiquetas fuera de elementos flotantes
% (etiquetas de figuras)
\usepackage{sidecap} % Para poner el texto de las imágenes al lado
	\sidecaptionvpos{figure}{c} % Para que el texto se alinie al centro vertical
\usepackage{caption} % Para poder quitar numeracion de figuras
\usepackage{commath} % funcionalidades extras para diferenciales, integrales,
% etc (\od, \dif, etc)
\usepackage{cancel} % para cancelar expresiones (\cancelto{0}{x})
 
\usepackage{anysize} 					% Para personalizar el ancho de  los márgenes
\marginsize{2cm}{2cm}{2cm}{2cm} % Izquierda, derecha, arriba, abajo

\usepackage{appendix}
\renewcommand{\appendixname}{Apéndices}
\renewcommand{\appendixtocname}{Apéndices}
\renewcommand{\appendixpagename}{Apéndices} 

% Para que las referencias sean hipervínculos a las figuras o ecuaciones y
% aparezcan en color
\usepackage[colorlinks=true,plainpages=true,citecolor=blue,linkcolor=blue]{hyperref}
%\usepackage{hyperref} 
% Para agregar encabezado y pie de página
\usepackage{fancyhdr} 
\pagestyle{fancy}
\fancyhf{}
\fancyhead[L]{\footnotesize UGR} %encabezado izquierda
\fancyhead[R]{\footnotesize ETSIIT}   % dereecha
\fancyfoot[R]{\footnotesize Memoria Prácticas}  % Pie derecha
\fancyfoot[C]{\thepage}  % centro
\fancyfoot[L]{\footnotesize Máster en Ing. Informática}  %izquierda
\renewcommand{\footrulewidth}{0.4pt}


\usepackage{listings} % Para usar código fuente
\definecolor{dkgreen}{rgb}{0,0.6,0} % Definimos colores para usar en el código
\definecolor{gray}{rgb}{0.5,0.5,0.5} 
% configuración para el lenguaje que queramos utilizar
\lstset{language=Matlab,
   keywords={break,case,catch,continue,else,elseif,end,for,function,
      global,if,otherwise,persistent,return,switch,try,while},
   basicstyle=\ttfamily,
   keywordstyle=\color{blue},
   commentstyle=\color{red},
   stringstyle=\color{dkgreen},
   numbers=left,
   numberstyle=\tiny\color{gray},
   stepnumber=1,
   numbersep=10pt,
   backgroundcolor=\color{white},
   tabsize=4,
   showspaces=false,
   showstringspaces=false}

\newcommand{\sen}{\operatorname{\sen}}	% Definimos el comando \sen para el seno
%en español

\title{Memoría para las prácticas de la asignatura Domótica}

%%%%%%%% TERMINA PREÁMBULO %%%%%%%%%%%%

\begin{document}

%%%%%%%%%%%%%%%%%%%%%%%%%%%%%%%%%% PORTADA %%%%%%%%%%%%%%%%%%%%%%%%%%%%%%%%%%%%%%%%%%%%
																					%%%
\begin{center}																		%%%
\newcommand{\HRule}{\rule{\linewidth}{0.5mm}}									%%%\left
 																					%%%
\begin{minipage}{0.48\textwidth} \begin{flushleft}
%\includegraphics[scale = 0.63]{Imagenes/logo_upiita}
\end{flushleft}\end{minipage}
\begin{minipage}{0.48\textwidth} \begin{flushright}

%\includegraphics[scale = 0.35]{Imagenes/IPN}
\end{flushright}\end{minipage}

													 								%%%
\vspace*{-1.5cm}								%%%
																					%%%	
\textsc{\huge Universidad de  \\ \vspace{5px} Granada}\\[1.5cm]	

\textsc{\LARGE Escuela Técnica Superior en Ingenier\'ias 				%%%
Inform\'atica \\ \vspace{5px} y de las Telecomunicaciones}\\[1.5cm]													%%%

\begin{minipage}{0.9\textwidth} 
\begin{center}																					%%%
\textsc{\LARGE Memoria Prácticas}
\end{center}
\end{minipage}\\[0.5cm]
%%%
    																				%%%
 			\vspace*{1cm}																		%%%
																					%%%
\HRule \\[0.4cm]																	%%%
{ \huge \bfseries Domótica}\\[0.4cm]	%%%
 																					%%%
\HRule \\[1.5cm]																	%%%
 																				%%%
																					%%%
\begin{minipage}{0.46\textwidth}													%%%
\begin{flushleft} \large															%%%
\emph{Autor:}\\	
Manuel Jesús García Manday\\
%%%
			%\vspace*{2cm}	
            													%%%
										 						%%%
\end{flushleft}																		%%%
\end{minipage}		
																%%%
\begin{minipage}{0.52\textwidth}		
\vspace{-0.6cm}											%%%
\begin{flushright} \large															%%%
%\emph{Asesores:} \\																	%%%
%Asesor 1\\
%Asesor 2	\\
%Asesor 3														%%%
\end{flushright}																	%%%
\end{minipage}	
\vspace*{1cm}
%\begin{flushleft}
 	
%\end{flushleft}
%%%
 		\flushleft{\textbf{\Large Máster en Ing. Informática}	}\\																		%%%
\vspace{2cm} 																				
\begin{center}																					
{\large \today}																	%%%
 			\end{center}												  						
\end{center}							 											
																					
\newpage																		
%%%%%%%%%%%%%%%%%%%% TERMINA PORTADA %%%%%%%%%%%%%%%%%%%%%%%%%%%%%%%%

\tableofcontents 

\newpage

\section{Objetivo.}

El objetivo de este documento es presentar una memoria sobre la realización de las prácticas de la asignatura. A través del mismo se describirán los pasos empleados para el desarrollo de cada una de las prácticas propuestas así como las herramientas utilizadas para ello. Para el desarrollo de las mismas se va a emplear el software ETS5 como entrono de desarrollo integrado (IDE) y el sistema KNX como entorno de pruebas y despliegue de los diferentes proyectos, con el fin de realizar la implementación y pruebas pertinentes para probar que el comportamiento obtenido por las distintas funcionalidades de las prácticas es el correcto. 



\section{Práctica 1. Regulación de luz.} 

La finalidad de esta primera práctica es la de realizar una funcionalidad que sea capaz de regular la intensidad luminosa de una lámpara led a través del módulo de botonera táctil

\section{Desarrollo de la solución.}

bla bla bla bla bla bla bla bla bla bla bla bla bla bla blabla bla bla bla blabla bla bla bla blabla bla bla bla blabla bla bla bla blabla bla bla bla blabla bla bla bla blabla bla bla bla blabla bla bla bla bla bla bla bla bla blabla bla bla bla bla bla bla bla bla blabla bla bla bla blabla bla bla bla blabla bla bla bla bla bla bla bla bla bla bla bla bla bla bla bla bla bla bla bla bla bla bla bla bla bla bla bla bla bla bla bla bla bla bla



\begin{figure}[H]
	\begin{center}
% 		\includegraphics[width = 0.5\textwidth]{Imagenes/IPN}
 		\captionof{figure}{\label{fig:IPN}Descripción de la imagen} 
	\end{center} 
\end{figure}



\subsection{Generación del concepto.}

bla bla bla bla bla bla bla bla bla bla bla bla bla bla blabla bla bla bla blabla bla bla bla blabla bla bla bla blabla bla bla bla blabla bla bla bla blabla bla bla bla blabla bla bla bla blabla bla bla bla bla bla bla bla bla blabla bla bla bla bla bla bla bla bla blabla bla bla bla blabla bla bla bla blabla bla bla bla bla bla bla bla bla bla bla bla bla bla bla bla bla bla bla bla bla bla bla bla bla bla bla bla bla bla bla bla bla bla bla

bla bla bla bla bla bla bla bla bla bla bla bla bla bla blabla bla bla bla blabla bla bla bla blabla bla bla bla blabla bla bla bla blabla bla bla bla blabla bla bla bla blabla bla bla bla blabla bla bla bla bla bla bla bla bla blabla bla bla bla bla bla bla bla bla blabla bla bla bla blabla bla bla bla blabla bla bla bla bla bla bla bla bla bla bla bla bla bla bla bla bla bla bla bla bla bla bla bla bla bla bla bla bla bla bla bla bla bla bla
bla bla bla bla bla bla bla bla bla bla bla bla bla bla blabla bla bla bla blabla bla bla bla blabla bla bla bla blabla bla bla bla blabla bla bla bla blabla bla bla bla blabla bla bla bla blabla bla bla bla bla bla bla bla bla blabla bla bla bla bla bla bla bla bla blabla bla bla bla blabla bla bla bla blabla bla bla bla bla bla bla bla bla bla bla bla bla bla bla bla bla bla bla bla bla bla bla bla bla bla bla bla bla bla bla bla bla bla bla
bla bla bla bla bla bla bla bla bla bla bla bla bla bla blabla bla bla bla blabla bla bla bla blabla bla bla bla blabla bla bla bla blabla bla bla bla blabla bla bla bla blabla bla bla bla blabla bla bla bla bla bla bla bla bla blabla bla bla bla bla bla bla bla bla blabla bla bla bla blabla bla bla bla blabla bla bla bla bla bla bla bla bla bla bla bla bla bla bla bla bla bla bla bla bla bla bla bla bla bla bla bla bla bla bla bla bla bla bla


bla bla bla bla bla bla bla bla bla bla bla bla bla bla blabla bla bla bla blabla bla bla bla blabla bla bla bla blabla bla bla bla blabla bla bla bla blabla bla bla bla blabla bla bla bla blabla bla bla bla bla bla bla bla bla blabla bla bla bla bla bla bla bla bla blabla bla bla bla blabla bla bla bla blabla bla bla bla bla bla bla bla bla bla bla bla bla bla bla bla bla bla bla bla bla bla bla bla bla bla bla bla bla bla bla bla bla bla bla


\section{Conclusiones.}

\section{Cronograma de trabajo terminal II.}


%%%%%%% Bibliografía %%%%%%%%
\bibliographystyle{bst/IEEEtran.bst} 
\addcontentsline{toc}{section}{Referencias}  
\bibliography{bib/IEEEabrv,bib/IEEEreferencias.bib} 
%%%%%%% Bibliografía %%%%%%%%    

\appendix  
\clearpage % o \cleardoublepage
\addappheadtotoc 
\appendixpage 

\section{Anexos 1.}


bla bla bla bla bla bla bla bla bla bla bla bla bla bla blabla bla bla bla blabla bla bla bla blabla bla bla bla blabla bla bla bla blabla bla bla bla blabla bla bla bla blabla bla bla bla blabla bla bla bla bla bla bla bla bla blabla bla bla bla bla bla bla bla bla blabla bla bla bla blabla bla bla bla blabla bla bla bla bla bla bla bla bla bla bla bla bla bla bla bla bla bla bla bla bla bla bla bla bla bla bla bla bla bla bla bla bla bla bla 

bla bla bla bla bla bla bla bla bla bla bla bla bla bla blabla bla bla bla blabla bla bla bla blabla bla bla bla blabla bla bla bla blabla bla bla bla blabla bla bla bla blabla bla bla bla blabla bla bla bla bla bla bla bla bla blabla bla bla bla bla bla bla bla bla blabla bla bla bla blabla bla bla bla blabla bla bla bla bla bla bla bla bla bla bla bla bla bla bla bla bla bla bla bla bla bla bla bla bla bla bla bla bla bla bla bla bla bla bla 

\section{Anexos 2.}


bla bla bla bla bla bla bla bla bla bla bla bla bla bla blabla bla bla bla blabla bla bla bla blabla bla bla bla blabla bla bla bla blabla bla bla bla blabla bla bla bla blabla bla bla bla blabla bla bla bla bla bla bla bla bla blabla bla bla bla bla bla bla bla bla blabla bla bla bla blabla bla bla bla blabla bla bla bla bla bla bla bla bla bla bla bla bla bla bla bla bla bla bla bla bla bla bla bla bla bla bla bla bla bla bla bla bla bla bla
bla bla bla bla bla bla bla bla bla bla bla bla bla bla blabla bla bla bla blabla bla bla bla blabla bla bla bla blabla bla bla bla blabla bla bla bla blabla bla bla bla blabla bla bla bla blabla bla bla bla bla bla bla bla bla blabla bla bla bla bla bla bla bla bla blabla bla bla bla blabla bla bla bla blabla bla bla bla bla bla bla bla bla bla bla bla bla bla bla bla bla bla bla bla bla bla bla bla bla bla bla bla bla bla bla bla bla bla bla
bla bla bla bla bla bla bla bla bla bla bla bla bla bla blabla bla bla bla blabla bla bla bla blabla bla bla bla blabla bla bla bla blabla bla bla bla blabla bla bla bla blabla bla bla bla blabla bla bla bla bla bla bla bla bla blabla bla bla bla bla bla bla bla bla blabla bla bla bla blabla bla bla bla blabla bla bla bla bla bla bla bla bla bla bla bla bla bla bla bla bla bla bla bla bla bla bla bla bla bla bla bla bla bla bla bla bla bla bla
bla bla bla bla bla bla bla bla bla bla bla bla bla bla blabla bla bla bla blabla bla bla bla blabla bla bla bla blabla bla bla bla blabla bla bla bla blabla bla bla bla blabla bla bla bla blabla bla bla bla bla bla bla bla bla blabla bla bla bla bla bla bla bla bla blabla bla bla bla blabla bla bla bla blabla bla bla bla bla bla bla bla bla bla bla bla bla bla bla bla bla bla bla bla bla bla bla bla bla bla bla bla bla bla bla bla bla bla bla


bla bla bla bla bla bla bla bla bla bla bla bla bla bla blabla bla bla bla blabla bla bla bla blabla bla bla bla blabla bla bla bla blabla bla bla bla blabla bla bla bla blabla bla bla bla blabla bla bla bla bla bla bla bla bla blabla bla bla bla bla bla bla bla bla blabla bla bla bla blabla bla bla bla blabla bla bla bla bla bla bla bla bla bla bla bla bla bla bla bla bla bla bla bla bla bla bla bla bla bla bla bla bla bla bla bla bla bla bla
bla bla bla bla bla bla bla bla bla bla bla bla bla bla blabla bla bla bla blabla bla bla bla blabla bla bla bla blabla bla bla bla blabla bla bla bla blabla bla bla bla blabla bla bla bla blabla bla bla bla bla bla bla bla bla blabla bla bla bla bla bla bla bla bla blabla bla bla bla blabla bla bla bla blabla bla bla bla bla bla bla bla bla bla bla bla bla bla bla bla bla bla bla bla bla bla bla bla bla bla bla bla bla bla bla bla bla bla bla

\end{document}